%%%%%%%%%%%%%%%%%%%%%%%%%%%%%%%%%%%%%%%%%%%%%%%%%%%%%%%%%%%%%%%%%%%%%%%%%%%%%%%
%
%   TWO-COLUMN EXAM CHEATSHEET TEMPLATE
%
%   This template is designed to maximize data density on an A4 page.
%   It uses small margins, no page numbers, and compact list/paragraph
%   settings.
%
%%%%%%%%%%%%%%%%%%%%%%%%%%%%%%%%%%%%%%%%%%%%%%%%%%%%%%%%%%%%%%%%%%%%%%%%%%%%%%%

% --- DOCUMENT CLASS ---
% 10pt: Small base font size
% a4paper: Standard A4 paper
% twocolumn: The core requirement
\documentclass[10pt, a4paper, twocolumn]{article}


% --- PACKAGE CONFIGURATION ---\

% 1. GEOMETRY: Set very small margins
% This is the most important setting for "max data".
% All margins (top, bottom, left, right) are set to 1cm.
\usepackage[a4paper, top=1cm, bottom=1cm, left=1cm, right=1cm]{geometry}
\usepackage{lipsum}

% 2. FONT & LANGUAGE (Universal Preamble Protocol)
% We use fontspec to load modern Noto fonts, ensuring compilation.
% \usepackage{fontspec}
\usepackage[english, bidi=basic, provide=*]{babel}
\babelprovide[import, onchar=ids fonts]{english}
% Use Noto Sans. It's very clear and compact (sans-serif).
% \babelfont{rm}{Noto Sans}
% \babelfont{sf}{Noto Sans}
% \babelfont{tt}{Noto Sans Mono}

% 3. TYPOGRAPHY & SPACING
% microtype improves justification and text spacing, making it look
% tighter and more professional.
\usepackage{microtype}
\usepackage{enumitem} % For list customization
\usepackage{amsmath} % For math
\usepackage{amssymb} % For math symbols
\usepackage{algorithmic} % For algorithms
\usepackage{algorithm} % For algorithm floats
\usepackage{booktabs} % For better tables
\usepackage{graphicx} % For images (if any)

% 4. THEOREM-LIKE ENVIRONMENTS (optional, but good for definitions)
\usepackage{amsthm}
\newtheoremstyle{crisp}
  {} % Space above
  {} % Space below
  {\itshape} % Body font
  {} % Indent amount
  {\bfseries} % Theorem head font
  {.} % Punctuation after theorem head
  {.5em} % Space after theorem head
  {} % Theorem head spec
\theoremstyle{crisp}
\newtheorem{theorem}{Theorem}
\newtheorem{definition}{Definition}
\newtheorem{remark}[theorem]{Remark}
\newtheorem{example}[theorem]{Example}
\newtheorem{conjecture}[theorem]{Conjecture}
\newtheorem{question}[theorem]{Question}

% 5. HYPERLINKS (Must be last)
% Useful if you want clickable links in your PDF (e.g., to web resources)
\usepackage{hyperref}
\hypersetup{
    colorlinks=true,
    linkcolor=blue,
    urlcolor=blue,
    pdftitle={Exam Cheatsheet},
    pdfauthor={Student},
}


% --- GLOBAL DOCUMENT SETTINGS ---

% 1. Page Style
% \pagestyle{empty} removes all headers and footers (no page numbers).
\pagestyle{empty}

% 2. List Spacing (CRITICAL FOR DENSITY)
% This removes *all* vertical spacing from itemize and enumerate.
\setlist{noitemsep, topsep=0pt, partopsep=0pt, parsep=0pt}

% 3. Paragraph Spacing
% \setlength{\parindent}{0pt} removes the first-line indent.
% \setlength{\parskip}{0.5ex} adds a tiny vertical space between paragraphs.
\setlength{\parindent}{0pt}
\setlength{\parskip}{0.5ex}

% 4. Column Spacing
% Adjusts the blank space *between* the two columns.
\setlength{\columnsep}{1.5em} % 1.5em is af good default


% --- DOCUMENT START ---
\begin{document}

% --- TABLE OF CONTENTS ---
\tableofcontents
\newpage

% --- MODULES START HERE ---

% --- MODULE 1: Algorithm Analysis & Amortization ---
% --- LECTURES 1, 2, 3, 9 & CLRS 16 ---

\section{Module 1: Algorithm Analysis \& Amortization}

\subsection{Core Concepts}
\begin{description}
    \item[Worst-Case Analysis] (Lec 2) Measures the maximum number of operations an algorithm will perform for any input of size $n$. This is the standard for analysis in this course.
    \item[Amortized Analysis] (Lec 9, CLRS 16) Analyzes a \textit{sequence} of operations to find the average cost per operation in the worst case. It is \textbf{not} the same as average-case analysis.
\end{description}

\subsection{Asymptotic Notations (Lec 3)}
\noindent For a given time function $T(n) \ge 0$:
\begin{description}
    \item[$O(f(n))$ (Big-Oh)] (Upper Bound) $\exists$ positive constants $c, n_0$ such that $T(n) \le c \cdot f(n)$ for all $n \ge n_0$.
    \item[$\Omega(f(n))$ (Big-Omega)] (Lower Bound) $\exists$ positive constants $c, n_0$ such that $T(n) \ge c \cdot f(n)$ for all $n \ge n_0$.
    \item[$\Theta(f(n))$ (Big-Theta)] (Tight Bound) $\exists$ positive constants $c_1, c_2, n_0$ s.t. $c_1 f(n) \le T(n) \le c_2 f(n)$ for all $n \ge n_0$. (This means $T(n)$ is $O(f(n))$ and $\Omega(f(n))$).
\end{description}

\subsection{Common Growth Rates (Lec 3)}
$O(1) \subset O(\log n) \subset O(n) \subset O(n \log n) \subset O(n^2) \subset O(n^k) \subset O(2^n) \subset O(n!)$

\subsection{Amortized Analysis Methods (Lec 9, CLRS 16)}
\subsubsection{1. Aggregate Analysis (Lec 9)}
Show that for a sequence of $n$ operations, the total cost $T(n)$ is $O(f(n))$. The amortized cost per operation is then $T(n)/n$.
\begin{itemize}
    \item \textbf{Binary Counter ($n$ increments):} Total cost = $\sum_{i=0}^{\log n} \frac{n}{2^i} \approx 2n = O(n)$. Amortized cost = $\mathbf{O(1)}$.
    \item \textbf{Dynamic Array (Push):} Total cost for $n$ pushes = $O(n)$ (for $n$ pushes) + $O(n)$ (for all copying). Amortized cost = $\mathbf{O(1)}$.
\end{itemize}

\subsubsection{2. Accounting Method (CLRS 16.2)}
Assign different \textit{amortized costs} ($\hat{c}$) to operations.
\begin{itemize}
    \item If $\hat{c} > c$ (actual cost), the extra charge is stored as \textbf{credit} on the data structure.
    \item If $\hat{c} < c$, the deficit is paid by stored credit.
    \item \textbf{Rule:} The total credit must \textbf{never be negative}.
    \item \textbf{Dynamic Array (Push):}
    \begin{itemize}
        \item Let \textbf{Amortized Cost $\hat{c} = 3$}.
        \item \textbf{Case 1: No resize (actual cost $c=1$)}. We pay 3. 1 pays for the push, 2 are stored as credit.
        \item \textbf{Case 2: Resize (actual cost $c=n+1$)}. Occurs when $n-1$ items are present. There must be $(n-1)$ items, each with 2 units of credit, for a total of $2(n-1)$ credit. We use $n$ of this credit to pay for the copy, and 1 to pay for the push.
        \item We use $n+1$ credit, but have $2n-2$. Credit remains non-negative. Amortized cost is $\mathbf{O(1)}$.
    \end{itemize}
\end{itemize}

\subsubsection{3. Potential Method (CLRS 16.3)}
Define a potential function $\Phi$ that maps the data structure $D$ to a number.
\begin{itemize}
    \item $\Phi(D_0) = 0$ and $\Phi(D_i) \ge 0$ for all $i$.
    \item Amortized cost: $\hat{c_i} = c_i + \Phi(D_i) - \Phi(D_{i-1})$ (actual cost + change in potential).
    \item \textbf{Dynamic Array (Push):}
    \begin{itemize}
        \item Define $\Phi(T) = 2 \cdot \text{num}(T) - \text{size}(T)$. (Assumes $\text{size}/2 < \text{num} \le \text{size}$).
        \item \textbf{Case 1: No resize ($c_i=1$)}. $\text{num}_i = \text{num}_{i-1}+1$, $\text{size}_i = \text{size}_{i-1}$.
        \item $\Delta\Phi = \Phi(D_i) - \Phi(D_{i-1}) = (2 \cdot \text{num}_i - \text{size}_i) - (2 \cdot (\text{num}_i-1) - \text{size}_i) = 2$.
        \item $\hat{c_i} = c_i + \Delta\Phi = 1 + 2 = 3$.
        \item \textbf{Case 2: Resize ($c_i=n$)}. $\text{num}_i = n$, $\text{size}_i = 2(n-1)$. $\text{num}_{i-1} = n-1$, $\text{size}_{i-1} = n-1$.
        \item $\Delta\Phi = (2n - 2(n-1)) - (2(n-1) - (n-1)) = 2 - (n-1) = 3 - n$.
        \item $\hat{c_i} = c_i + \Delta\Phi = n + (3 - n) = 3$.
        \item In both cases, the amortized cost is $\mathbf{O(1)}$.
    \end{itemize}
\end{itemize}

% --- MODULE 2: Linear Data Structures ---
% --- LECTURES 4, 5, 6, 7, 8, 10 & CLRS 10 ---

\section{Module 2: Linear Data Structures (Lists, Stacks, Queues)}

\subsection{Arrays vs. Linked Lists (Lec 4, 5, 6)}
A fundamental trade-off in data structures.
\begin{center}
\begin{tabular}{|l|l|l|}
\hline
\textbf{Operation} & \textbf{Array} & \textbf{Linked List} \\ \hline
Read (by index) & $O(1)$ & $O(n)$ \\
Search & $O(n)$ & $O(n)$ \\
Insert (at front) & $O(n)$ & $O(1)$ \\
Insert (at end) & $O(1)$* & $O(1)$** \\
Delete (at front) & $O(n)$ & $O(1)$ \\
Delete (at end) & $O(1)$ & $O(n)$*** \\
\hline
\end{tabular}
\end{center}
\begin{footnotesize}
*Amortized $O(1)$ with dynamic array (Lec 9).
**Requires a \texttt{tail} pointer.
***Requires traversal to find predecessor, or a Doubly Linked List.
\end{footnotesize}
\begin{itemize}
    \item \textbf{Sentinel Nodes (CLRS 10.2):} To simplify code, use a special dummy node `L.nil` for the `head` and to mark the `next` of the tail. This removes edge cases for empty lists or single-node lists.
\end{itemize}

\subsection{Stack ADT (LIFO - Last-In, First-Out) (Lec 6, 7)}
\begin{description}
    \item[Push(e)] Insert element at the top. \textbf{$O(1)$*}.
    \item[Pop()] Remove and return element from top. \textbf{$O(1)$}.
    \item[Top()] Return top element without removing. \textbf{$O(1)$}.
\end{description}
\begin{footnotesize}
*Amortized $O(1)$ using a dynamic array (Lec 9).
\end{footnotesize}

\subsubsection{Application: Balanced Parentheses (Lec 7)}
Algorithm to check if a string of brackets is balanced.
\begin{itemize}
    \item Create an empty stack.
    \item Iterate through the string:
    \begin{itemize}
        \item If char is an \textbf{opening} bracket ('(', '[', '\{'), \textbf{push} it onto the stack.
        \item If char is a \textbf{closing} bracket (')', ']', '\}'):
        \begin{itemize}
            \item If stack is empty $\to$ \textbf{return false}.
            \item \textbf{Pop} the top element.
            \item If popped bracket does not match the closing bracket $\to$ \textbf{return false}.
        \end{itemize}
    \end{itemize}
    \item After the loop, \textbf{return true} if the stack is empty, else \textbf{false}.
\end{itemize}
\textbf{Complexity:} $O(n)$ time, $O(n)$ space.

\subsection{Queue ADT (FIFO - First-In, First-Out) (Lec 8)}
\begin{description}
    \item[Enqueue(e)] Insert element at the back. \textbf{$O(1)$*}.
    \item[Dequeue()] Remove and return element from front. \textbf{$O(1)$**}.
    \item[Front()] Return front element without removing. \textbf{$O(1)$}.
\end{description}
\begin{footnotesize}
*Amortized $O(1)$ using a dynamic/doubling array.
**Requires a \textbf{circular array} implementation to avoid an $O(n)$ shift.
\end{footnotesize}

\subsubsection{Application: Maze Solving / BFS (Lec 10)}
\begin{itemize}
    \item Queues are the core data structure for \textbf{Breadth-First Search (BFS)}.
    \item BFS explores "layer by layer", which is used to find the \textbf{shortest path} in unweighted graphs (like a maze).
    \item \textbf{Proof of Correctness:} This is a non-trivial proof by induction. See the proof in \textbf{Module 7} (Graphs).
\end{itemize}

\subsection{General Tree Representation (CLRS 10.3)}
How to store a tree where nodes have an \textit{arbitrary} number of children?
\begin{itemize}
    \item \textbf{Bad Idea:} Array of pointers at each node (wastes space).
    \item \textbf{Good Idea: Left-Child, Right-Sibling Representation.}
    \item Each node $x$ has only two pointers:
    \begin{enumerate}
        \item $x.\text{left\_child}$: points to its first child.
        \item $x.\text{right\_sibling}$: points to its next sibling.
    \end{enumerate}
    \item This creates a binary tree, where all "left" pointers go down one level, and all "right" pointers stay on the same level.
\end{itemize}

% --- MODULE 3: Trees, Heaps, & Priority Queues ---
% --- LECTURES 12, 13, 14, 15, 16 & CLRS 6 ---

\section*{Module 3: Trees, Heaps, \& Priority Queues}

\subsection*{Tree ADT (Lec 12, 13)}
A data structure that stores elements hierarchically.
\begin{description}
    \item[Root] The single top-most node.
    \item[Leaf (External)] A node with no children.
    \item[Internal Node] A node with one or more children.
    \item[Depth of $v$] Length of the path from root to $v$. `depth(root) = 0`.
    \item[Height of $v$] Length of the longest path from $v$ to a leaf. `height(leaf) = 0`.
    \item[Height of Tree] Height of the root.
\end{description}

\subsection*{Binary Trees (Lec 13, 14)}
A tree where each node has at most two children: a \textbf{left} child and a \textbf{right} child.
\begin{itemize}
    \item \textbf{Property:} A binary tree of height $h$ has $n$ nodes, where
    $h+1 \le n \le 2^{h+1} - 1$.
    \item This implies the height $h$ is at least $\Omega(\log n)$ and at most $O(n)$.
\end{itemize}

\subsection*{Tree Traversals (Lec 13, 14)}
Systematic ways to visit all nodes in a tree.
\begin{description}
    \item[Preorder] (Root, Left, Right). Used to copy a tree.
    \item[Inorder] (Left, Root, Right). Used in BSTs to get a sorted list.
    \item[Postorder] (Left, Right, Root). Used to delete a tree or calculate directory sizes (Lec 13).
    \item[Level-Order] (BFS). Visits nodes layer by layer using a \textbf{Queue}.
\end{description}

\subsection*{Binary Heap (Min-Heap) (Lec 14-16, CLRS 6)}
A \textbf{Priority Queue} is often implemented using a Min-Heap.
\begin{description}
    \item[Structure Property] It is an \textbf{Almost Complete Binary Tree (ACBT)}. All levels are full except possibly the last, which is filled from left to right.
    \item[Height] An ACBT with $n$ nodes has height $h = \lfloor \log n \rfloor$. (CLRS 6.1)
    \item[Order Property (Min-Heap)] For every node $u$ (except the root), $\text{key}(\text{parent}(u)) \le \text{key}(u)$. The minimum element is always at the root.
\end{description}

\subsubsection*{Array-Based Representation (Lec 15)}
An ACBT can be stored efficiently in an array. For a node at 0-based index $i$:
\begin{itemize}
    \item $\text{parent}(i) = \lfloor (i-1) / 2 \rfloor$
    \item $\text{left}(i) = 2i + 1$
    \item $\text{right}(i) = 2i + 2$
\end{itemize}

\subsubsection*{Heap Operations (Lec 15, 16, CLRS 6.2-6.3)}
\begin{tabular}{l l l}
    \textbf{Operation} & \textbf{Algorithm} & \textbf{Time} \\
    \hline
    \texttt{GetMin} & Read root & $O(1)$ \\
    \texttt{Insert(k)} & Add to end, \textbf{Trickle-Up} & $O(\log n)$ \\
    \texttt{DeleteMin} & Replace root with last, & $O(\log n)$ \\
     & \textbf{Trickle-Down (Heapify)} & \\
    \texttt{BuildHeap} & Call `Heapify` from & $\mathbf{O(n)}$ \\
     & $n/2$ down to 0 & \\
    \hline
\end{tabular}

\subsubsection*{Correctness of Heap Operations}
\begin{itemize}
    \item \textbf{Trickle-Up (Insert):} After inserting $k$, the only violation is if $k < \text{parent}(k)$. Swapping $k$ with its parent fixes this violation and moves the \textit{only} potential violation one level up. This repeats until $k \ge \text{parent}(k)$ or $k$ is the root.
    
    \item \textbf{Heapify(A, i) (Trickle-Down, CLRS 6.2):}
    \begin{itemize}
        \item \textbf{Assumption:} The subtrees at `left(i)` and `right(i)` are already valid heaps.
        \item \textbf{Goal:} Make the tree at $i$ a valid heap.
        \item \textbf{Logic:} Find the smallest of $\{i, \text{left}(i), \text{right}(i)\}$. If $i$ is not the smallest, swap $A[i]$ with $A[\text{smallest}]$. This fixes the violation at $i$.
        \item \textbf{Proof:} The swap \textit{might} violate the heap property at the `smallest` child's new position. By recursively calling `Heapify(A, smallest)`, we correct this. This continues down the tree until a leaf is reached.
        \item \textbf{Time:} $O(h) = O(\log n)$.
    \end{itemize}
    
    \item \textbf{Build-Heap(A) (Lec 16, CLRS 6.3):}
    \begin{itemize}
        \item \textbf{Algorithm:} Call `Heapify(i)` for $i = \lfloor n/2 \rfloor$ down to $0$.
        \item \textbf{Correctness (Loop Invariant):} "At the start of each iteration $i$, nodes $i+1, \dots, n-1$ are all roots of valid max-heaps."
        \item \textbf{Init:} True. All nodes from $\lfloor n/2 \rfloor + 1$ to $n-1$ are leaves, which are trivial heaps.
        \item \textbf{Maintenance:} `Heapify(i)` assumes children at $2i+1$ and $2i+2$ are heap roots. Since $2i+1 > i$, the invariant holds. `Heapify` makes $i$ a heap root, maintaining the invariant for the next (smaller) $i$.
        \item \textbf{Termination:} When $i=0$, node 0 is made a heap root. Since $1 \dots n-1$ are already heap roots, the entire tree is a valid heap.
        \item \textbf{Time Analysis:} $\mathbf{O(n)}$, not $O(n \log n)$. Most nodes are near the bottom and have small height. $\sum_{h=0}^{\log n} (n/2^{h+1}) \cdot O(h) = O(n)$.
    \end{itemize}
\end{itemize}

% --- MODULE 4: Heap Applications & Greedy Algorithms ---
% --- LECTURES 17, 18, 19, 20 & CLRS 6.4, 15.3 ---

\section*{Module 4: Heap Apps \& Greedy Algorithms}

\subsection*{Heap Sort (Lec 17, CLRS 6.4)}
An in-place sorting algorithm with $O(n \log n)$ worst-case time.

\subsubsection*{Algorithm (for ascending order)}
\begin{enumerate}
    \item \textbf{Build Max-Heap:} Convert the input array $A$ of $n$ elements into a max-heap. ($O(n)$)
    \item \textbf{Iterate $i = (n-1)$ down to $1$:}
    \begin{itemize}
        \item \textbf{Swap} root ($A[0]$, the max element) with $A[i]$.
        \item \textbf{Shrink} the heap size by 1 (max element is now sorted at $A[i]$).
        \item \textbf{Heapify(A, 0)} on the new root to restore the max-heap property on the $A[0 \dots i-1]$ subarray. ($O(\log n)$)
    \end{itemize}
\end{enumerate}
\textbf{Total Time:} $O(n) \text{ (Build)} + (n-1) \times O(\log n) \text{ (Heapify)} = \mathbf{O(n \log n)}$.

\subsubsection*{Proof of Correctness}
\textbf{Loop Invariant:} At the start of each iteration $i$:
\begin{enumerate}
    \item The subarray $A[i+1 \dots n-1]$ contains the $(n-i)$ largest elements of $A$, \textit{in sorted order}.
    \item The subarray $A[0 \dots i]$ contains the $i+1$ smallest elements and is a valid max-heap.
\end{enumerate}
\begin{itemize}
    \item \textbf{Init:} $i=n-1$. The invariant is trivially true (subarray $A[n \dots n-1]$ is empty). `Build-Heap` made $A[0 \dots n-1]$ a heap.
    \item \textbf{Maintenance:} Swapping $A[0]$ (the max) with $A[i]$ puts the next largest element in its correct sorted spot. `Heapify(0)` restores the heap property for $A[0 \dots i-1]$.
    \item \textbf{Termination:} When $i=0$, the loop terminates. The array $A[1 \dots n-1]$ contains the $n-1$ largest elements in sorted order. $A[0]$ must be the smallest.
\end{itemize}


\subsection*{Greedy Algorithm: Huffman Coding (Lec 17-20, CLRS 15.3)}
A \textbf{greedy algorithm} that generates an optimal prefix-free binary code for data compression.
\begin{description}
    \item[Problem] Given symbols with frequencies $p_i$, find a binary code for each symbol to minimize the average encoding length: $\min \sum p_i \cdot (\text{depth of symbol } i)$.
    \item[Prefix-Free Codes] No code is a prefix of another (e.g., if `01` is a code, `011` cannot be).
    \item[Tree Representation] Prefix-free codes correspond to \textbf{full binary trees} where symbols are \textbf{only at the leaves}. The code for a symbol is the path from the root (0=left, 1=right).
\end{description}

\subsubsection*{Huffman's Algorithm (Greedy Strategy)}
Builds the optimal tree from the bottom up.
\begin{enumerate}
    \item \textbf{Initialize:} Create a \textbf{min-priority queue} (Min-Heap) of $n$ "trees", where each tree is a single leaf node for a symbol, and its priority is its frequency.
    \item \textbf{Iterate $n-1$ times:}
    \begin{itemize}
        \item \texttt{DeleteMin} twice to get the two trees ($T_x, T_y$) with the \textbf{lowest frequencies} ($p_x, p_y$).
        \item \textbf{Merge} them: Create a new parent node with $T_x$ and $T_y$ as children.
        \item The new tree's frequency is $p_x + p_y$.
        \item \textbf{Insert} this new, merged tree back into the min-priority queue.
    \end{itemize}
    \item \textbf{Finish:} The last tree remaining in the queue is the optimal Huffman tree.
\end{enumerate}
\textbf{Running Time:} $O(n \log n)$, because it involves $O(n)$ `Insert` and $O(n)$ `DeleteMin` operations on a priority queue of size $O(n)$.

\subsubsection*{Proof of Correctness (CLRS 15.3)}
Uses induction on $n=|\Sigma|$. Proves two key properties:
\begin{description}
    \item[Greedy Choice Property] In *some* optimal tree, the two lowest-frequency symbols ($x, y$) are sibling leaves at the maximum depth.
    \begin{itemize}
        \item \textit{Proof (by exchange):} Take any optimal tree $T^*$. Let $b, c$ be two siblings at max depth. Let $x, y$ be the two lowest-frequency symbols. If $\{b,c\} \neq \{x,y\}$, we know $p_x \le p_b$ and $p_y \le p_c$. Swap $x$ with $b$ and $y$ with $c$. The new cost can only decrease or stay the same. Thus, an optimal tree with $x, y$ as deep siblings exists.
    \end{itemize}
    \item[Optimal Substructure] Let $T'$ be an optimal tree for the subproblem where $x,y$ are replaced by a new symbol $z$ with $p_z = p_x + p_y$. Let $T$ be the tree formed by replacing $z$ in $T'$ with a node that has children $x, y$.
    \begin{itemize}
        \item We have $L(T) = L(T') + p_x + p_y$.
        \item The cost of $T$ is the cost of $T'$ plus a \textit{fixed constant} ($p_x+p_y$).
        \item Therefore, minimizing $L(T)$ is equivalent to minimizing $L(T')$.
    \end{itemize}
    \item[Conclusion] The algorithm makes the (safe) greedy choice and then optimally solves the resulting subproblem (by induction). Thus, the final tree $T$ is optimal.
\end{description}

% --- MODULE 5: Hash Tables ---
% --- LECTURES 21, 22, 23 ---

\section{Module 5: Hash Tables}

\subsection{Core Idea \& Operations}
The main goal is to support `Insert`, `Delete`, and `Lookup` (Search) in average-case $\mathbf{O(1)}$ time.
\begin{itemize}
    \item \textbf{Setup:} We have a large \textbf{Universe $U$} of possible keys. We want to store a (much smaller) dynamic set $S \subseteq U$.
    \item \textbf{Structure:} An array $A$ (buckets) of size $n$, where $n \approx |S|$.
    \item \textbf{Hash Function:} A function $h: U \to \{0, \dots, n-1\}$ that maps keys to array indices (buckets).
\end{itemize}

\subsection{Applications}
\begin{description}
    \item[De-duplication (Lec 21)] Check if a new item in a stream has been seen before.
    \begin{itemize}
        \item For new item $x$, `Lookup(x)`. If not found, `Insert(x)`.
        \item $\mathbf{O(1)}$ per item.
    \end{itemize}
    \item[2-SUM Problem (Lec 21)] Given array $A$ (size $n$) and target $t$, find if $x, y \in A$ s.t. $x+y=t$.
    \begin{itemize}
        \item `Insert` all elements of $A$ into a hash table $H$. ($O(n)$)
        \item For each $x \in A$, `Lookup(t-x)` in $H$. ($O(n) \times O(1) = O(n)$)
        \item \textbf{Total Time: $\mathbf{O(n)}$}. (Beats $O(n \log n)$ sorting approach).
    \end{itemize}
\end{description}

\subsection{Collisions (Lec 21, 23)}
Collisions are inevitable (by Pigeonhole Principle) when $|U| > n$. The \textbf{Birthday Paradox} shows they are far more likely than expected (e.g., $n=10k$, 400 objects $\to$ 99.9\% collision chance).

\subsection{Collision Resolution Strategies (Lec 22)}

\subsubsection{1. Chaining (Separate Chaining)}
\begin{itemize}
    \item Each bucket $A[i]$ holds a pointer to a \textbf{linked list} of all elements that hash to $i$.
    \item \textbf{Load Factor:} $\alpha = |S| / n$. Can be $> 1$.
    \item \textbf{Operations:}
    \begin{itemize}
        \item `Insert(x)`: Insert $x$ at the head of list $A[h(x)]$. \textbf{Time: $O(1)$}.
        \item `Lookup(x)`: Search list $A[h(x)]$.
        \item `Delete(x)`: Search and delete from list $A[h(x)]$.
    \end{itemize}
    \item \textbf{Performance:} All ops depend on list length. If $h$ is good ("spreads data"), average list length is $\alpha$.
    \item \textbf{Expected Time:} $\mathbf{O(1+\alpha)}$. If we maintain $n = \Theta(|S|)$, then $\alpha=O(1)$, and all ops are \textbf{$O(1)$}.
\end{itemize}

\subsubsection{2. Open Addressing}
\begin{itemize}
    \item At most one object per bucket. Requires $n \ge |S|$ (so $\alpha \le 1$).
    \item \textbf{Probing:} If $h(x)$ is full, probe for the next empty slot using a probe sequence.
    \item \textbf{Linear Probing:} Try $h(x), (h(x)+1)\%n, (h(x)+2)\%n, \dots$
    \item \textbf{Double Hashing:} Try $h_1(x), (h_1(x)+h_2(x))\%n, (h_1(x)+2h_2(x))\%n, \dots$
    \item \textbf{Delete:} Tricky. Cannot just empty the slot. Must mark it as "DELETED".
    \item \textbf{Performance:} Heavily degrades as $\alpha \to 1$. Needs $\alpha \ll 1$.
\end{itemize}

\subsection{Good Hash Functions (Lec 23)}
\begin{description}
    \item[Goal] ``Spreads out the data'' evenly, mapping similar keys to different buckets.
    
    \item[Bad Functions] 
    \texttt{h(x) = first 3 digits} (bad for phone numbers) or 
    \texttt{h(x) = x mod n} where $n = 2^k$ (bad for data with common low-bit patterns).
    
    \item[Method] \texttt{h(x) = (hash\_code(x)) mod n}
    \begin{itemize}
        \item \textbf{Hash Code:} Map key $x$ (e.g., string) to a large integer.
        \item \textbf{Compression:} Map the integer to $\{0, \dots, n-1\}$.
    \end{itemize}
    
    \item[Choosing $n$ (bucket count)] 
    $n$ should be a \textbf{prime number} that is not too close to a power of 2.  
    This helps break up patterns in the input data.
    
    \item[Universal Hashing] 
    To defeat a ``pathological dataset’’ (adversary), we use a \textit{family} of hash functions $\mathcal{H}$ and pick one $h \in \mathcal{H}$ at random.  
    This guarantees good average performance, no matter the input data.
\end{description}


% --- MODULE 6: Binary Search Trees, AVL Trees, B-Trees ---
% --- LECTURES 24-32 & CLRS 12, 18 & GTM 10.4 ---

\section{Module 6: Search Trees (BST, AVL, B-Trees)}

\subsection{Binary Search Tree (BST) (Lec 24, 25, CLRS 12)}
A binary tree that maintains a sorted order.

\subsubsection{The BST Invariant}
For any node $x$ in the tree:
\begin{itemize}
    \item All keys in $x$'s \textbf{left} subtree are $\le \text{key}(x)$.
    \item All keys in $x$'s \textbf{right} subtree are $\ge \text{key}(x)$.
\end{itemize}
\textbf{Property (CLRS Thm 12.1):} An \texttt{InorderTraversal} of a BST outputs its elements in sorted order.
\textit{Proof (by induction):} A subtree of 1 node is sorted. Assume Inorder($T_L$) and Inorder($T_R$) are sorted. Then Inorder($T$) = [Inorder($T_L$)] + [root] + [Inorder($T_R$)]. By the BST invariant, all items in $T_L \le$ root $\le$ all items in $T_R$. Thus, the full list is sorted.

\subsubsection{Operations (all $O(h)$, where $h$ is tree height)}
\begin{description}
    \item[\texttt{Search(k)}] Start at root. If $k < \text{key}(x)$, go left. If $k > \text{key}(x)$, go right. If $x=$ NULL, $k$ is not found.
    \item[\texttt{Minimum(x)}] Start at $x$, follow \texttt{left} child pointers until NULL.
    \item[\texttt{Maximum(x)}] Start at $x$, follow \texttt{right} child pointers until NULL.
    \item[\texttt{Insert(z)}] \texttt{Search} for key $k$. The (NULL) leaf position where the search ends is where $z$ is inserted.
\end{description}
\textbf{Theorem 12.4 (CLRS):} The expected height of a \textbf{randomly built} BST on $n$ distinct keys is $\mathbf{O(\log n)}$. This means that on average, all operations are efficient.

\subsubsection{Successor(x) (Lec 25, CLRS 12.2)}
Finds the node with the smallest key that is \textit{larger} than $\text{key}(x)$.
\begin{itemize}
    \item \textbf{Case 1: $x$ has a right subtree.}
    The successor is the \texttt{Minimum()} node in $x$'s right subtree.
    \item \textbf{Case 2: $x$ has NO right subtree.}
    Go up the tree using \texttt{parent} pointers. The successor is the \textbf{first ancestor} $y$ such that $x$ is in the \textbf{left subtree} of $y$. (This is the first "right turn" you take on the path up). If you reach the root and no such $y$ exists, $x$ was the maximum.
\end{itemize}

\subsubsection{Delete(z) (Lec 25, CLRS 12.3)}
This is the most complex operation, formalized by CLRS with a helper `TRANSPLANT(u, v)` which replaces subtree $u$ with subtree $v$.
\begin{itemize}
    \item \textbf{Case 1: $z$ has no left child.}
    Replace $z$ with its right child (which could be NULL).
    \item \textbf{Case 2: $z$ has no right child.}
    Replace $z$ with its left child.
    \item \textbf{Case 3: $z$ has two children.}
    \begin{enumerate}
        \item Find $z$'s \textbf{successor}, $y = \texttt{Minimum}(z.\text{right})$.
        \item $y$ is guaranteed to be in $z$'s right subtree and will have \textbf{no left child}.
        \item \textbf{If $y$ is $z$'s direct child:} Replace $z$ with $y$, and $y$ adopts $z$'s left child.
        \item \textbf{If $y$ is not $z$'s direct child:}
        \begin{itemize}
            \item Replace $y$ with its own right child (a Case 1 op).
            \item $y$ adopts $z$'s right child.
            \item Replace $z$ with $y$, and $y$ adopts $z$'s left child.
        \end{itemize}
    \end{enumerate}
\end{itemize}
\textbf{Problem:} All operations are $O(h)$. If keys are inserted in sorted order, the tree becomes a linked list with $h=O(n)$. Operations degrade to $O(n)$.

\subsection{Augmented BSTs (Lec 26)}
\textbf{Idea:} Store extra data in each node to answer new types of queries.
\begin{itemize}
    \item \textbf{Rule:} The augmented data must be computable from the node's key and its children's augmented data in $O(1)$ time.
    \item \textbf{Example: Subtree Size.} Store $size(x)$ at each node.
    \item \textbf{Update Rule:} $size(x) = size(\text{left}(x)) + size(\text{right}(x)) + 1$.
    \item \textbf{New Operation: \texttt{Rank(t)}} (How many nodes are $\le t$?) in $O(h)$.
\end{itemize}

\subsection{AVL Trees (Self-Balancing BST) (Lec 27, 28, 29)}
\textbf{Goal:} Guarantee $h = O(\log n)$ by keeping the tree balanced.
\begin{description}
    \item[AVL Property (Height-Balance)] For \textit{every} node $x$, the heights of its children can differ by at most 1:
    $|\text{height}(\text{left}(x)) - \text{height}(\text{right}(x))| \le 1$.
    \item[Complexity] All BST operations are guaranteed $\mathbf{O(\log n)}$.
\end{description}

\subsubsection{Insert \& Delete in AVL Trees}
\begin{enumerate}
    \item Perform the standard BST `Insert` or `Delete`.
    \item Walk back up from the modified node to the root, updating the `height` of each ancestor.
    \item Find the \textbf{first (lowest) ancestor $z$} that violates the AVL property.
    \item Perform the correct \textbf{Rotation} operation at $z$ to restore balance.
    \item \textit{Proof of Correctness (Insert):} `Insert` requires at most 1 (single or double) rotation. The rotation at $z$ is designed such that the \textbf{new height} of the rotated subtree is the \textbf{same as the height of $z$ before the insertion}. This means no ancestors above $z$ become unbalanced.
    \item \textit{Delete:} May require $O(\log n)$ rotations as you travel up, as a rotation at one node might unbalance its parent.
\end{enumerate}
\subsubsection{Rotations (The Fix)}
Let $z$ be the unbalanced node, $y$ be its taller child, and $x$ be $y$'s taller child.
\begin{description}
    \item[Case 1: Left-Left] (Single \textbf{Right Rotation} at $z$)
    \item[Case 2: Right-Right] (Single \textbf{Left Rotation} at $z$)
    \item[Case 3: Left-Right] (\textbf{Left Rotation} at $y$, then \textbf{Right Rotation} at $z$)
    \item[Case 4: Right-Left] (\textbf{Right Rotation} at $y$, then \textbf{Left Rotation} at $z$)
\end{description}

\subsection{(2,4) Trees (Lec 30, 31, GTM 10.4)}
A balanced M-way tree ($M=4$) where nodes can hold multiple keys.
\begin{description}
    \item[Size Property] Every internal node has 2, 3, or 4 children (and 1, 2, or 3 keys).
    \item[Depth Property] All leaves are at the \textbf{same level}.
    \item[Height] $h = \mathbf{O(\log n)}$ (proven by $\min$ 2 children/node $\to$ $h \le \log_2 n$).
\end{description}

\subsubsection{Insert in (2,4)-Tree}
\begin{enumerate}
    \item \texttt{Search} to find the correct leaf node to insert $k$.
    \item \textbf{If leaf has < 3 keys:} Add $k$ to the node. Done.
    \item \textbf{If leaf is full (3 keys):} This is an \textbf{Overflow}.
    \begin{enumerate}
        \item Add $k$ to the full node (temporarily a "5-node" with 4 keys).
        \item \textbf{Split} the node: Promote the middle key ($k_2$) up to the parent.
        \item The remaining keys form two new nodes (with $k_1$ and $k_3, k_4$) which are attached to the parent.
        \item If the parent is now full, \textbf{propagate the split upwards}.
        \item If the root splits, the tree height grows by 1.
    \end{enumerate}
\end{enumerate}
\textbf{Correctness:} The split operation is the only structural change, and it perfectly maintains the (2,4) properties (size and depth).

\subsubsection{Delete (Bottom-Up) (GTM 10.4.3)}
\begin{enumerate}
    \item \texttt{Search} for key $k$. If $k$ is not in a leaf, swap it with its successor (which must be in a leaf).
    \item Remove $k$ from the leaf node.
    \item \textbf{If leaf is now empty (0 keys):} This is an \textbf{Underflow}.
    \begin{itemize}
        \item \textbf{Case 1: Transfer.} If an adjacent sibling has $\ge 2$ keys, perform a "rotation". Move a key from the parent down to the empty node, and move a key from the sibling up to the parent.
        \item \textbf{Case 2: Fusion.} If adjacent siblings are minimal (only 1 key), \textbf{merge} the empty node, its minimal sibling, and the key from the parent.
        \item This fusion may cause the parent to underflow, so the process may \textbf{propagate upwards}. If the root underflows and is empty, it is deleted and the height shrinks by 1.
    \end{itemize}
\end{enumerate}

\subsubsection{Delete (Top-Down / Pre-emptive)}
\textbf{Idea:} Ensures we never delete from a minimal node (1-key node), avoiding upward propagation.
\begin{enumerate}
    \item \textbf{Handle Non-Leaf:} If $k$ is in an internal node, swap it with its inorder successor (which is in a leaf). The problem is now to delete the successor from a leaf.
    \item \textbf{Pre-emptive Traversal:} Start at the root. As we traverse \textit{down} to the target leaf, at each node $x$ we visit:
    \item \textbf{If $x$ is minimal (1 key):}
    \begin{itemize}
        \item \textbf{Case 1 (Transfer):} If an adjacent sibling has $\ge 2$ keys, perform a transfer (rotation) through the parent.
        \item \textbf{Case 2 (Fusion):} If adjacent siblings are also minimal, perform a fusion with a sibling and the parent's separating key. (If the root is part of a fusion and becomes empty, the fused child becomes the new root.)
    \end{itemize}
    \item Now $x$ is guaranteed to have $\ge 2$ keys.
    \item \textbf{Descend:} Move to the appropriate child (which may have changed after a fusion).
    \item \textbf{Final Deletion:} When we reach the target leaf, it is \textit{guaranteed} to have $\ge 2$ keys. Remove $k$. No underflow occurs.
\end{enumerate}


\subsection{B-Trees (Lec 32, CLRS 18)}
The generalization of (2,4)-trees, ideal for disk-based memory.
\begin{description}
    \item[Properties (min-degree $t$)]
    \begin{itemize}
        \item Each node holds $k$ keys, where $t-1 \le k \le 2t-1$.
        \item Each internal node with $k$ keys has $k+1$ children.
        \item All leaves are at the same depth.
        \item Root has at least 1 key (or 0 if tree is empty).
    \end{itemize}
    \item[Height] $h = O(\log_t n)$. This is excellent for disk, as $t$ can be large (e.g., 1000), making $h$ very small.
\end{description}

\subsubsection{Search in B-Tree (CLRS 18.2)}
$O(t)$ work per node, but $O(1)$ disk access. Total disk accesses: $\mathbf{O(\log_t n)}$.

\subsubsection{Insert in B-Tree (CLRS 18.3)}
\textbf{Idea:} Avoids the "underflow-propagation" of (2,4) trees by being \textbf{pre-emptive}. We never insert into a full node.
\begin{enumerate}
    \item Start at root. As we traverse \textit{down} to find the leaf:
    \item \textbf{If we encounter a full node $x$ ($2t-1$ keys):}
    \item \textbf{Split $x$ first} using `B-TREE-SPLIT-CHILD`.
    \item `B-TREE-SPLIT-CHILD(parent, i, x)`:
    \begin{itemize}
        \item Splits $x$ (the $i$-th child of parent) into two nodes, $x$ and $z$.
        \item $x$ gets the first $t-1$ keys. $z$ gets the last $t-1$ keys.
        \item The median key of $x$ (key $t$) moves up into the parent.
    \end{itemize}
    \item After splitting, we continue our search from the (now non-full) parent.
    \item \textbf{Finally:} We reach the target leaf, which is \textit{guaranteed} not to be full. Insert key $k$.
\end{enumerate}
\textbf{Correctness:} This maintains all B-Tree properties. The only way height grows is if the root itself splits.

% --- MODULE 7: Graph Algorithms (Expanded) ---
% --- LECTURES 33-41 & CLRS 19-22 ---

\section*{Module 7: Graph Algorithms}

\subsection*{Graph Definitions (Lec 33)}
\begin{description}
    \item[Graph] $G = (V, E)$. $n = |V|, m = |E|$.
    \item[Path] Sequence of vertices $(v_1, \dots, v_k)$ where $(v_i, v_{i+1}) \in E$.
    \item[Cycle] A path that starts and ends at the same vertex. A graph with no cycles is \textbf{acyclic}. A \textbf{DAG} is a Directed Acyclic Graph.
    \item[Handshake Lemma] $\sum_{v \in V} \text{deg}(v) = 2m$.
\end{description}

\subsection*{Graph Representations (Lec 34)}
\begin{description}
    \item[Adjacency List] Array of $n$ linked lists. $Adj[u]$ stores neighbors of $u$.
    \begin{itemize}
        \item \textbf{Memory:} $\mathbf{O(n+m)}$. \textbf{Standard for sparse graphs.}
        \item \textbf{Check Edge $(u, v)$:} $O(\text{deg}(u))$.
    \end{itemize}
    \item[Adjacency Matrix] $n \times n$ matrix $A$ where $A[u][v] = 1$ if $(u, v) \in E$.
    \begin{itemize}
        \item \textbf{Memory:} $\mathbf{O(n^2)}$. \textbf{Standard for dense graphs.}
        \item \textbf{Check Edge $(u, v)$:} $O(1)$.
    \end{itemize}
\end{description}

\subsection*{Breadth-First Search (BFS) (Lec 35, CLRS 20.2)}
\textbf{Idea:} Explores in "layers," 1 edge away, then 2, etc.
\begin{itemize}
    \item \textbf{Data Structure:} A \textbf{FIFO Queue}.
    \item \textbf{Algorithm (from $s$):}
    \begin{enumerate}
        \item Mark $s$ as explored (GRAY), `dist(s)=0`, `enqueue(s)`.
        \item While queue is not empty:
        \item $u \gets \texttt{dequeue()}$.
        \item For each neighbor $v$ of $u$:
        \item If $v$ is unexplored (WHITE):
        \item Mark $v$ GRAY, `dist(v) = dist(u) + 1`, `parent(v)=u`, `enqueue(v)`.
        \item Mark $u$ as finished (BLACK).
    \end{enumerate}
    \item \textbf{Complexity:} $\mathbf{O(n+m)}$.
\end{itemize}
\textbf{Application: Shortest Paths (Unweighted).}
\begin{itemize}
    \item \textbf{Correctness (White Path Theorem):} BFS finds the shortest path $d(s,v)$ to all reachable $v$.
    \item \textbf{Proof (by Induction on distance $k$):} (Lec 11/12)
    \item \textbf{Hypothesis $Hyp(k)$:} BFS finds all nodes at distance $k$ (set $S_k$) after finding all nodes at distance $k-1$ (set $S_{k-1}$), and `dist(v)=k` $\forall v \in S_k$.
    \item \textbf{Base $k=1$:} $S_0=\{s\}$ is enqueued. Its neighbors ($S_1$) are enqueued next, all with `dist = 0+1 = 1`. True.
    \item \textbf{Step $Hyp(k) \to Hyp(k+1)$:} The queue ensures all nodes from $S_k$ are processed before any from $S_{k+1}$. When a $u \in S_k$ is dequeued, it discovers all its unvisited neighbors $v$. These $v$ must be in $S_{k+1}$ (if they were in $S_k$ or earlier, they'd be visited). BFS sets `dist(v) = dist(u) + 1 = k+1`. True.
\end{itemize}

\subsection*{Depth-First Search (DFS) (Lec 37, CLRS 20.3)}
\textbf{Idea:} Explores as far as possible before backtracking.
\begin{itemize}
    \item \textbf{Data Structure:} A \textbf{LIFO Stack} (or recursion).
    \item \textbf{Algorithm (recursive, from $s$):}
    \begin{enumerate}
        \item Mark $s$ as explored (GRAY) and set discovery time $d[s]$.
        \item For each neighbor $v$ of $s$:
        \item If $v$ is unexplored (WHITE):
        \item `parent(v)=u`, call `DFS(v)`.
        \item Mark $s$ as finished (BLACK) and set finish time $f[s]$.
    \end{enumerate}
    \item \textbf{Complexity:} $\mathbf{O(n+m)}$.
\end{itemize}
\textbf{Properties (CLRS 20.3):}
\begin{itemize}
    \item \textbf{Parenthesis Theorem:} The discovery/finish intervals are nested. For any two nodes $u, v$, either their intervals are disjoint, or one is contained within the other (iff one is a descendant of the other in the DFS tree).
    \item \textbf{Edge Types:}
    \begin{itemize}
        \item \textbf{Tree Edge:} To a WHITE child.
        \item \textbf{Back Edge:} To a GRAY ancestor (indicates a \textbf{CYCLE}).
        \item \textbf{Forward/Cross Edge:} To a BLACK node.
    \end{itemize}
\end{itemize}

\subsection*{DFS Applications (Lec 38, 39, CLRS 20.4, 20.5)}
\begin{description}
    \item[Topological Sort (for DAGs)] A linear ordering where for all edges $(u, v)$, $u$ appears before $v$.
    \begin{itemize}
        \item \textbf{Algorithm:} Run DFS. As each vertex finishes, insert it onto the \textbf{front} of a linked list.
        \item \textbf{Correctness:} For any edge $(u,v)$, DFS must finish $v$ before it finishes $u$ (if it saw $(u,v)$ and $v$ was GRAY, it's a cycle). Thus, $f[v] < f[u]$. Since the list is sorted by decreasing $f[]$, $u$ comes before $v$.
    \end{itemize}
    \item[Strongly Connected Components (SCCs)] (Kosaraju's)
    \begin{enumerate}
        \item Run DFS on $G$ to get finish times $f[u]$.
        \item Compute $G^T$ (reverse all edges).
        \item Run DFS on $G^T$, visiting vertices in order of \textbf{decreasing} $f[u]$.
        \item Each tree in the 2nd DFS's forest is one SCC.
    \end{enumerate}
    \item \textbf{Correctness:} (Lec 39) The 1st DFS finds a "source" SCC (the one with the highest finish time $u$). The 2nd DFS (on $G^T$) starts at $u$ and can only explore $u$'s SCC, because all edges now point *into* that component from other SCCs, so the DFS gets "trapped" in the SCC.
\end{description}

\subsection*{Weighted Shortest Paths: Dijkstra (CLRS 22.3)}
Finds shortest paths from $s$ in a \textbf{weighted} graph with \textbf{no negative edges}.
\begin{itemize}
    \item \textbf{Data Structure:} A \textbf{Min-Priority Queue} (Min-Heap).
    \item \textbf{Algorithm (from $s$):}
    \begin{enumerate}
        \item Init: `dist[v] = $\infty$` for all $v$, `dist[s] = 0`.
        \item Add all $v$ to Min-Heap $Q$ (keyed by `dist[v]`).
        \item While $Q$ is not empty:
        \item $u \gets Q.\texttt{DeleteMin()}$.
        \item For each neighbor $v$ of $u$:
        \item \textbf{Relaxation:} If `dist[u] + weight(u,v) < dist[v]`:
        \item `dist[v] = dist[u] + weight(u,v)`
        \item `parent(v) = u`
        \item \texttt{DecreaseKey}$(Q, v)$
    \end{enumerate}
    \item \textbf{Complexity:} $O(n \log n + m \log n)$ (naive heap) or $\mathbf{O(n \log n + m)}$ (with Fibonacci heap or as $m \ge n$ $\mathbf{O(m \log n)}$ with binary heap).
    \item \textbf{Correctness (Cut Property):} When Dijkstra adds $u$ to the set of "finished" nodes, `dist[u]` is the true shortest path. (Proof: Assume not. Let $u$ be the first node added with the wrong distance. Consider the true shortest path $s \dots x \to y \dots u$, where $x$ is finished and $y$ is not. Because $x$ is finished, $d(s,y) \le d(s,u)$. And since $d(s,y) < d(s,u)$ (non-neg edges), the algorithm would have chosen $y$ before $u$. Contradiction.)
\end{itemize}
\textbf{Special Case: 0-1 BFS} : If all weights are 0 or 1, use a \textbf{Deque}. When relaxing an edge: if weight is 0, `push\_front(v)`; if weight is 1, `push\_back(v)`. This maintains the "layers" and achieves $O(n+m)$.

\subsection*{Minimum Spanning Tree (MST) (Lec 40, CLRS 21)}
Finds a min-cost spanning tree in a connected, undirected, weighted graph.
\begin{itemize}
    \item \textbf{Generic Correctness (Cut Property):} For any cut (partition of $V$ into $S, V-S$), if $(u,v)$ is the \textbf{unique} minimum-weight edge crossing the cut, it \textbf{must} be in the MST. (Proof: Assume not. Adding $(u,v)$ to the MST $T^*$ creates a cycle. This cycle must have another edge $(x,y)$ crossing the cut. $c(x,y) > c(u,v)$. Swapping $(u,v)$ for $(x,y)$ gives a cheaper tree. Contradiction.)
\end{itemize}

\subsubsection*{Kruskal's Algorithm (Lec 40, 41, CLRS 21.2)}
\textbf{Idea:} A greedy algorithm that builds the MST by adding the cheapest edges that don't form a cycle.
\begin{enumerate}
    \item \textbf{Sort} all $m$ edges by weight.
    \item Init: $T = \emptyset$. Init a \textbf{DSU} with $n$ sets (one per vertex).
    \item For each edge $(u, v)$ in sorted order:
    \item If $\texttt{Find}(u) \neq \texttt{Find}(v)$:
    \item $T = T \cup \{(u, v)\}$
    \item $\texttt{Union}(u, v)$
\end{enumerate}
\textbf{Time:} $\mathbf{O(m \log m)}$ or $\mathbf{O(m \log n)}$ (dominated by edge sort).

\subsubsection*{Prim's Algorithm (CLRS 21.2)}
\textbf{Idea:} A greedy algorithm that "grows" the MST from a single vertex $s$.
\begin{enumerate}
    \item \textbf{Data Structure:} A \textbf{Min-Priority Queue} (Min-Heap) $Q$.
    \item Init: For all $v$, `key[v] = $\infty$`, `parent[v] = NULL`. Set `key[s] = 0`.
    \item Add all $v$ to $Q$ (keyed by `key[v]`).
    \item While $Q$ is not empty:
    \item $u \gets Q.\texttt{DeleteMin()}$.
    \item For each neighbor $v$ of $u$:
    \item If $v \in Q$ and $\text{weight}(u,v) < \text{key}[v]$:
    \item `parent[v] = u`, `key[v] = weight(u,v)`
    \item \texttt{DecreaseKey}$(Q, v)$
\end{enumerate}
\textbf{Time:} $\mathbf{O(m \log n)}$ (with a binary heap).

\subsection*{Disjoint Set Union (DSU) (Lec 41, CLRS 19)}
Data structure for maintaining disjoint sets. Used by Kruskal's.
\begin{description}
    \item[\texttt{MakeSet(x)}] Create a new set $\{x\}$.
    \item[\texttt{Find(x)}] Return the representative (root) of $x$'s set.
    \item[\texttt{Union(x, y)}] Merge the sets containing $x$ and $y$.
\end{description}
\textbf{Optimizations (Required for good performance):}
\begin{enumerate}
    \item \textbf{Union-by-Size (or Rank):} In `Union(x, y)`, make the root of the \textbf{smaller} tree point to the root of the \textbf{larger} tree.
    \begin{itemize}
        \item \textbf{Proof:} This guarantees tree height is $O(\log n)$. A node's depth only increases when its set is merged into a larger one, at least doubling the set size. This can only happen $O(\log n)$ times.
    \end{itemize}
    \item \textbf{Path Compression:} During `Find(x)`, make every node on the path from $x$ to the root point directly to the root.
\end{enumerate}
\textbf{Complexity:} With both optimizations, a sequence of $m$ operations (on $n$ sets) takes $\mathbf{O(m \cdot \alpha(n))}$ time, where $\alpha(n)$ is the inverse Ackermann function, which is practically $O(1)$ (grows *extremely* slowly, e.g., $\alpha(n) \le 4$ for any practical $n$).

% --- MODULE 8 (PART 1): Problems (Arrays, Lists, Stacks, Heaps, Hashes, DP) ---
% --- This file is now UPDATED with problems from GFG PDFs ---

\section*{Module 8 (Part 1): Problems (Pre-BST \& DP)}

\subsection*{Module 1/2: Arrays, Search, Stacks, \& Lists}

\subsubsection*{Problem: 2-Sum (Lec 21, Blind 75)}
\begin{itemize}
\item \textbf{Problem:} Given an array of integers nums and a target, return indices of two numbers such that they add up to target. (Also called "Check for Pair with Given Sum" in GFG).
\item \textbf{Solution Idea:} Use a Hash Map (Module 5). Iterate through the array once. For each element x, check if target - x is already in the map. If yes, return the indices. If no, add x and its index to the map.
\item \textbf{Complexity:} $O(n)$ time, $O(n)$ space.
\end{itemize}

\subsubsection*{Problem: 3Sum (Blind 75)}
\begin{itemize}
\item \textbf{Problem:} Find all unique triplets in nums that sum to zero.
\item \textbf{Solution Idea:} Sort the array first ($O(n \log n)$). Then, iterate through the array with a main pointer i. For each i, use two new pointers j = i+1 and k = n-1. Move j and k toward each other, checking the sum. If sum < 0, move j right. If sum > 0, move k left.
\item \textbf{Key:} Must skip duplicate elements for i, j, and k to ensure unique triplets.
\item \textbf{Complexity:} $O(n^2)$ time, $O(\log n)$ or $O(n)$ space for sorting.
\end{itemize}

\subsubsection*{Problem: Search in Rotated Sorted Array (Blind 75)}
\begin{itemize}
\item \textbf{Problem:} Given a sorted array rotated at some pivot, search for a target in $O(\log n)$ time.
\item \textbf{Solution Idea:} Use a modified Binary Search (Module 1). In each step, check if the mid element is in the "left" sorted portion or the "right" sorted portion.
\item \begin{enumerate}
\item Find mid. If nums[mid] == target, return.
\item \textbf{If left half is sorted} (nums[low] <= nums[mid]):
\begin{itemize}
\item If target is in this half (nums[low] <= target < nums[mid]), search left (high = mid - 1).
\item Otherwise, search right (low = mid + 1).
\end{itemize}
\item \textbf{If right half is sorted} (nums[mid] < nums[high]):
\begin{itemize}
\item If target is in this half (nums[mid] < target <= nums[high]), search right (low = mid + 1).
\item Otherwise, search left (high = mid - 1).
\end{itemize}
\end{enumerate}
\item \textbf{Complexity:} $O(\log n)$ time, $O(1)$ space.
\end{itemize}

\subsubsection*{Problem: Find Majority Element (from GFG Arrays.pdf)}
\begin{itemize}
\item \textbf{Problem:} Find the element that appears more than $n/2$ times.
\item \textbf{Solution Idea (Boyer-Moore):} Maintain a candidate and a count. Iterate the array. If count == 0, set candidate = nums[i]. If nums[i] == candidate, count++. Else, count--.
\item \textbf{Key:} A second pass is needed to verify if the candidate truly appears > $n/2$ times, (though not if one is guaranteed to exist).
\item \textbf{Complexity:} $O(n)$ time, $O(1)$ space.
\end{itemize}

\subsubsection*{Problem: K-th Smallest of Two Sorted Arrays (Hard)}
\begin{itemize}
\item \textbf{Problem:} Given two sorted arrays $A$ (size $n$) and $B$ (size $m$), find the $k$-th smallest element in their union.
\item \textbf{Solution Idea:} Use Binary Search on the \textit{smaller} array (assume $A$, $n \le m$) to find a partition $i$.
\item We binary search for $i$ in the range $[\max(0, k-m), \min(k, n)]$.
\item Let $j = k - i$. We have found the correct partition when:
$A[i-1] \le B[j]$ AND $B[j-1] \le A[i]$.
\item (Handle edge cases $i=0, i=n, j=0, j=m$ with $-\infty$ and $+\infty$).
\item When found, the answer is $\max(A[i-1], B[j-1])$.
\item \textbf{Complexity:} $\mathbf{O(\log(\min(n, m)))}$ time, $O(1)$ space.
\end{itemize}

\subsubsection*{Problem: Container With Most Water (Blind 75)}
\begin{itemize}
\item \textbf{Problem:} Find two lines in height array which, with the x-axis, form a container holding the most water.
\item \textbf{Solution Idea:} Use a "two-pointer" approach. Start with l=0 and r=n-1. Calculate the area. To maximize area, you must move the pointer corresponding to the \textbf{shorter} line.
\item \textbf{Complexity:} $O(n)$ time, $O(1)$ space.
\end{itemize}

\subsubsection*{Problem: Trapping Rain Water (Blind 75, GFG Arrays.pdf)}
\begin{itemize}
\item \textbf{Problem:} Given an elevation map, compute how much water it can trap.
\item \textbf{Solution Idea (Optimal):} Use two pointers l=0 and r=n-1, and two variables max\_l and max\_r. If max\_l < max\_r, process the left pointer l (water is max\_l - height[l]) and move l right. Otherwise, process and move r left.
\item \textbf{Complexity:} $O(n)$ time, $O(1)$ space.
\end{itemize}

\subsubsection*{Problem: Product of Array Except Self (Blind 75, GFG Arrays.pdf)}
\begin{itemize}
\item \textbf{Problem:} Given nums, return answer where answer[i] is the product of all elements of nums except nums[i]. Must be $O(n)$ and no division.
\item \textbf{Solution Idea:} Use two passes. (1) Left-to-right pass to store prefix products in answer. (2) Right-to-left pass, multiplying by a running suffix product.
\item \textbf{Complexity:} $O(n)$ time, $O(1)$ space (if output array doesn't count).
\end{itemize}

\subsubsection*{Problem: Count Inversions (from GFG Divide-And-Conquer.pdf)}
\begin{itemize}
\item \textbf{Problem:} Find the number of pairs $(i, j)$ such that $i < j$ and $A[i] > A[j]$.
\item \textbf{Solution Idea:} Use a modified Merge Sort. During the merge step, when an element from the \textbf{right} subarray (R[j]) is moved before an element from the \textbf{left} subarray (L[i]), it means R[j] is smaller than all remaining elements in the left subarray.
\item if L[i] <= R[j]: move L[i].
\item else: move R[j], inversions += (mid - i + 1).
\item \textbf{Complexity:} $O(n \log n)$ time (same as Merge Sort).
\end{itemize}

\subsubsection*{Problem: Set Matrix Zeroes (LeetCode)}
\begin{itemize}
\item \textbf{Problem:} If an element in an $M \times N$ matrix is 0, set its entire row and column to 0. Do this in-place.
\item \textbf{Solution Idea:} Use the first row and first column of the matrix itself as storage to mark rows/cols for zeroing. Use two boolean flags to handle whether the first row/col themselves originally contained a zero.
\item \textbf{Complexity:} $O(M \cdot N)$ time, $O(1)$ space.
\end{itemize}

\subsubsection*{Problem: Valid Parentheses (Lec 7, Blind 75, GFG Stack.pdf)}
\begin{itemize}
\item \textbf{Problem:} Check if a string of '(', ')', '[', ']', '{', '}' is valid.
\item \textbf{Solution Idea:} This is the canonical Stack application. Iterate the string. Push opening brackets. When a closing bracket is found, pop and check for a match. If stack is empty on pop, or no match, or stack not empty at end, return false.
\item \textbf{Complexity:} $O(n)$ time, $O(n)$ space.
\end{itemize}

\subsubsection*{Problem: Next Greater Element (from GFG Stack.pdf)}
\begin{itemize}
\item \textbf{Problem:} For each element in an array, find the first element to its right that is greater than it.
\item \textbf{Solution Idea:} Use a Stack. Iterate the array from \textbf{right to left}.
\item For each element x:
\item while stack is not empty and stack.top() <= x: pop() from stack.
\item if stack is empty: result[i] = -1.
\item else: result[i] = stack.top().
\item push(x) onto the stack.
\item \textbf{Complexity:} $O(n)$ time (each element is pushed/popped once), $O(n)$ space.
\end{itemize}

\subsubsection*{Problem: Largest Rectangular Area in Histogram (from GFG Stack.pdf)}
\begin{itemize}
\item \textbf{Problem:} Find the largest rectangle of all-1s in a binary matrix (or largest area in a histogram).
\item \textbf{Solution Idea:} (For Histogram) Use a stack to store indices. We want to find the nearest smaller bar on the \textit{left} and \textit{right} for each bar i.
\item Iterate i from 0 to $n$: while stack not empty and hist[stack.top()] >= hist[i]:
\item tp = stack.pop(). The hist[tp] bar's right bound is i, left bound is stack.top().
\item area = hist[tp] * (i - stack.top() - 1). Update max\_area.
\item push(i) onto stack.
\item \textbf{Complexity:} $O(n)$ time, $O(n)$ space.
\end{itemize}

\subsubsection*{Problem: Evaluate Reverse Polish Notation (LeetCode)}
\begin{itemize}
\item \textbf{Problem:} Evaluate an arithmetic expression in Reverse Polish Notation (RPN).
\item \textbf{Solution Idea:} Classic Stack problem. Iterate tokens. If number, push. If operator, pop two numbers, perform operation, push result.
\item \textbf{Complexity:} $O(n)$ time, $O(n)$ space.
\end{itemize}

\subsubsection*{Problem: Linked List Cycle (Lec 37, Blind 75, GFG Linked-List.pdf)}
\begin{itemize}
\item \textbf{Problem:} Determine if a linked list has a cycle.
\item \textbf{Solution Idea:} Floyd's Tortoise and Hare. slow moves 1 step, fast moves 2. If they meet, a cycle exists.
\item \textbf{Complexity:} $O(n)$ time, $O(1)$ space.
\end{itemize}

\subsubsection*{Problem: Linked List Cycle II (LeetCode)}
\begin{itemize}
\item \textbf{Problem:} If a cycle exists, return the node where it begins.
\item \textbf{Solution Idea:} Find the meeting point meet. Reset a new pointer start to head. Move start and meet one step at a time. Their new meeting point is the cycle's start.
\item \textbf{Complexity:} $O(n)$ time, $O(1)$ space.
\end{itemize}

\subsubsection*{Problem: Reverse a Linked List (from GFG Linked-List.pdf)}
\begin{itemize}
\item \textbf{Problem:} Reverse a singly linked list.
\item \textbf{Solution Idea (Iterative):} Use three pointers: prev = NULL, curr = head, next = NULL.
\item while (curr != NULL): next = curr.next, curr.next = prev, prev = curr, curr = next.
\item Return prev (which is the new head).
\item \textbf{Complexity:} $O(n)$ time, $O(1)$ space.
\end{itemize}

\subsubsection*{Problem: Merge Two Sorted Lists (from GFG Linked-List.pdf)}
\begin{itemize}
\item \textbf{Problem:} Merge two sorted linked lists into one sorted list.
\item \textbf{Solution Idea:} Create a dummy head node. Maintain a tail pointer. Compare l1.val and l2.val. Append the smaller node to tail, and advance that list's pointer.
\item \textbf{Complexity:} $O(n+m)$ time, $O(1)$ space.
\end{itemize}

\subsubsection*{Problem: Intersection Point of Two Lists (from GFG Linked-List.pdf)}
\begin{itemize}
\item \textbf{Problem:} Find the node at which two linked lists intersect.
\item \textbf{Solution Idea 1 (Hash):} Add all nodes of l1 to a HashSet. Iterate l2 and return the first node found in the set. $O(n+m)$ time, $O(n)$ space.
\item \textbf{Solution Idea 2 (Pointers):} Get lengths len1, len2. Move the pointer of the longer list by abs(len1-len2) steps. Now, move both pointers one step at a time. The node where they meet is the intersection. $O(n+m)$ time, $O(1)$ space.
\end{itemize}

\subsubsection*{Problem: Remove Nth Node From End of List (Blind 75)}
\begin{itemize}
\item \textbf{Problem:} Remove the $n$-th node from the end of a list in one pass.
\item \textbf{Solution Idea:} Use two pointers. Move fast $n$ steps ahead. Then move fast and slow together until fast.next is NULL. slow is now at the node before the target.
\item \textbf{Complexity:} $O(L)$ time (one pass), $O(1)$ space.
\end{itemize}

\subsubsection*{Problem: Implement Queue using Stacks (from GFG Queue.pdf)}
\begin{itemize}
\item \textbf{Problem:} Implement Enqueue and Dequeue using two stacks.
\item \textbf{Solution Idea (Amortized $O(1)$ Dequeue):} Use s1 (input) and s2 (output).
\item \textbf{Enqueue(x):} s1.push(x). $O(1)$.
\item \textbf{Dequeue():} if s2 is empty: while s1 is not empty: s2.push(s1.pop()).
\item return s2.pop(). $O(1)$ amortized.
\item \textbf{Complexity:} Amortized $O(1)$ for both operations.
\end{itemize}

\subsection*{Module 3/4: Heaps \& Priority Queues}

\subsubsection*{Problem: Find Median from Data Stream (Blind 75, GFG Heaps.pdf)}
\begin{itemize}
\item \textbf{Problem:} Design a class to support addNum(int) and findMedian() from a stream.
\item \textbf{Solution Idea:} Use two heaps: a \textbf{Max-Heap} left\_half (smaller half) and a \textbf{Min-Heap} right\_half (larger half). Rebalance after each add so sizes differ by at most 1.
\item \textbf{Complexity:} addNum is $O(\log n)$, findMedian is $O(1)$.
\end{itemize}

\subsubsection*{Problem: Merge k Sorted Lists (Blind 75, GFG Heaps.pdf)}
\begin{itemize}
\item \textbf{Problem:} Merge $k$ sorted linked lists into one.
\item \textbf{Solution Idea:} Use a \textbf{Min-Priority Queue} (Min-Heap) of size $k$. Insert the head of all $k$ lists. While heap is not empty: DeleteMin node $u$, add $u$ to result, Insert $u.next$ into heap.
\item \textbf{Complexity:} $O(N \log k)$ time, where $N$ is total elements. $O(k)$ space.
\end{itemize}

\subsubsection*{Problem: Top K Frequent Elements (Blind 75)}
\begin{itemize}
\item \textbf{Problem:} Given nums and $k$, return the $k$ most frequent elements.
\item \textbf{Solution Idea (Heap):} (1) Build a Hash Map of (number, frequency) ($O(n)$). (2) Create a \textbf{Min-Heap} of size $k$. (3) Iterate the map, pushing (freq, num). If heap.size() > k, pop().
\item \textbf{Complexity:} $O(n \log k)$ time, $O(n+k)$ space.
\end{itemize}

\subsubsection*{Problem: Connect n Ropes with Minimum Cost (from GFG Heaps.pdf)}
\begin{itemize}
\item \textbf{Problem:} Given $n$ ropes of different lengths, connect them into one rope with minimum cost (cost of connecting $a$ and $b$ is $a+b$).
\item \textbf{Solution Idea:} This is Huffman's idea (Lec 19). Use a \textbf{Min-Priority Queue}.
\item Insert all $n$ rope lengths into the heap.
\item while heap.size() > 1: a = heap.DeleteMin(), b = heap.DeleteMin(). new\_rope = a + b. total\_cost += new\_rope. heap.Insert(new\_rope).
\item \textbf{Complexity:} $O(n \log n)$ time.
\end{itemize}

\subsection*{Module 5: Hash Tables}

\subsubsection*{Problem: Longest Substring Without Repeating Characters (Blind 75)}
\begin{itemize}
\item \textbf{Problem:} Find the length of the longest substring without repeating characters.
\item \textbf{Solution Idea:} Use the "Sliding Window" technique with a Hash Set (or Hash Map char -> index). Maintain window $[l, r]$. Move r. If char[r] is a duplicate (in set or map), move l to the right past the first occurrence.
\item \textbf{Complexity:} $O(n)$ time, $O(k)$ space (where $k$ is size of charset).
\end{itemize}

\subsubsection*{Problem: LRU Cache (Hard, Blind 75, GFG Hashing.pdf)}
\begin{itemize}
\item \textbf{Problem:} Design a cache with get and put in $O(1)$ time. It must evict the "Least Recently Used" item when full.
\item \textbf{Solution Idea:} Combine a \textbf{Hash Map} (map<key, Node*>) for $O(1)$ lookup and a \textbf{Doubly Linked List (DLL)} to maintain recency order.
\item \textbf{get(key):} Look up node in map. Move node to head of DLL.
\item \textbf{put(key, val):} If key exists, update value, move node to head. If new: create node, add to head. If full: evict tail of DLL, remove tail.key from map.
\item \textbf{Complexity:} $O(1)$ for both get and put.
\end{itemize}

\subsubsection*{Problem: Find Duplicates in O(n) time (from GFG Hashing.pdf)}
\begin{itemize}
\item \textbf{Problem:} Given an array of $n$ numbers (from 1 to $n$), find all duplicates.
\item \textbf{Solution Idea (In-place Hash):} Iterate the array. For each element x, go to index idx = abs(x) - 1. If nums[idx] is negative, x is a duplicate. Otherwise, mark nums[idx] as negative (nums[idx] = -nums[idx]).
\item \textbf{Complexity:} $O(n)$ time, $O(1)$ space.
\end{itemize}

\subsection*{Module X: Dynamic Programming (New Section)}

\subsubsection*{Problem: Coin Change (Blind 75)}
\begin{itemize}
\item \textbf{Problem:} Find the fewest number of coins to make a given amount, given coin denominations.
\item \textbf{Solution Idea (Bottom-Up DP):} Create dp[amount+1] initialized to $\infty$. Set dp[0] = 0.
\item for i from 1 to amount:
\item for c in coins:
\item if i - c >= 0:
\item dp[i] = min(dp[i], 1 + dp[i-c])
\item The answer is dp[amount].
\item \textbf{Complexity:} $O(A \cdot C)$ time, $O(A)$ space, where $A$=amount, $C$=num coins.
\end{itemize}

\subsubsection*{Problem: Longest Increasing Subsequence (LIS) (Blind 75)}
\begin{itemize}
\item \textbf{Problem:} Find the length of the longest strictly increasing subsequence.
\item \textbf{Solution Idea 1 (DP):} dp[i] = length of LIS \textit{ending at index i}. $O(n^2)$.
\item \textbf{Solution Idea 2 (Patience Sorting):} Maintain an array tails storing the smallest ending element for all LIS of a given length. For each x in nums, find the first tails[i] >= x and replace it (using binary search). If no such element, append x.
\item \textbf{Complexity 2:} $O(n \log n)$ time, $O(n)$ space.
\end{itemize}

\subsubsection*{Problem: Longest Common Subsequence (LCS) (Blind 75)}
\begin{itemize}
\item \textbf{Problem:} Find the length of the LCS between text1 and text2.
\item \textbf{Solution Idea (DP):} dp[i][j] = LCS of text1[0..i] and text2[0..j].
\item if text1[i] == text2[j]: dp[i][j] = 1 + dp[i-1][j-1]
\item else: dp[i][j] = max(dp[i-1][j], dp[i][j-1])
\item \textbf{Complexity:} $O(n \cdot m)$ time, $O(n \cdot m)$ space.
\end{itemize}

\subsubsection*{Problem: Word Break (Blind 75)}
\begin{itemize}
\item \textbf{Problem:} Can s be segmented into a sequence of dictionary words?
\item \textbf{Solution Idea (DP):} dp[i] = boolean, true if s[0..i-1] can be broken.
\item dp[0] = true. for i from 1 to n: for j from 0 to i-1:
\item if dp[j] and s[j..i-1] in wordDict: dp[i] = true; break;
\item \textbf{Complexity:} $O(n^2)$ time (due to substring check), $O(n)$ space.
\end{itemize}
% --- MODULE 8 (PART 2): Problems (BST, Tries, Graphs) ---
% --- This file is now UPDATED with problems from GFG PDFs ---

\section{Module 8 (Part 2): Problems (Post-BST)}

\subsection{Module 6: Trees, BSTs, Tries}

\subsubsection{Problem: Validate Binary Search Tree (Blind 75)}
\begin{itemize}
\item \textbf{Problem:} Determine if a binary tree is a valid BST.
\item \textbf{Solution Idea:} Use a recursive DFS. Pass down the valid range (min\_val, max\_val) allowed for the current node. The invariant must hold for all ancestors.
\item isValid(node, min, max):
\item if node.key <= min || node.key >= max: return false.
\item return isValid(node.left, min, node.key) \&\& isValid(node.right, node.key, max).
\item \textbf{Complexity:} $O(n)$ time, $O(h)$ space.
\end{itemize}

\subsubsection{Problem: BST Deletion (from GFG Binary-Search-Tree.pdf)}
\begin{itemize}
\item \textbf{Problem:} Delete a node z from a BST (Lec 25).
\item \textbf{Solution Idea:}
\begin{itemize}
\item \textbf{Case 1 (0 or 1 child):} Replace z with its child.
\item \textbf{Case 2 (2 children):} Find z's \textbf{inorder successor} y (the minimum node in z.right). y is guaranteed to have at most one (right) child.
\item Copy y.key into z.key.
\item Recursively delete y (which is now a Case 1 problem) from the right subtree.
\end{itemize}
\item \textbf{Complexity:} $O(h)$ time.
\end{itemize}

\subsubsection{Problem: Kth Smallest Element in a BST (Blind 75)}
\begin{itemize}
\item \textbf{Problem:} Find the $k$-th smallest element in a BST.
\item \textbf{Solution Idea 1 (Inorder):} Perform an inorder traversal. Stop at the $k$-th visited node. $O(n)$ time.
\item \textbf{Solution Idea 2 (Augmented BST):} (Lec 26) If nodes store size(left\_subtree), you can find it in $O(h)$ time.
\item \textbf{Complexity 2:} $O(h)$ time.
\end{itemize}

\subsubsection{Problem: Lowest Common Ancestor of a BST (from GFG)}
\begin{itemize}
\item \textbf{Problem:} Find the LCA of two nodes p and q in a BST.
\item \textbf{Solution Idea:} This is much simpler than a general BT. Start from the root.
\item if p.key < root.key \&\& q.key < root.key: recurse on root.left.
\item if p.key > root.key \&\& q.key > root.key: recurse on root.right.
\item else: root is the LCA (the point where p and q split).
\item \textbf{Complexity:} $O(h)$ time, $O(1)$ space (iterative).
\end{itemize}

\subsubsection{Problem: Convert Sorted Array to Balanced BST (from GFG)}
\begin{itemize}
\item \textbf{Problem:} Given a sorted array, create a height-balanced BST.
\item \textbf{Solution Idea:} This is a Divide and Conquer approach.
\begin{enumerate}
\item The \textbf{middle} element of the array arr[mid] becomes the root.
\item The left subarray arr[0...mid-1] recursively forms the root.left subtree.
\item The right subarray arr[mid+1...n-1] recursively forms the root.right subtree.
\end{enumerate}
\item \textbf{Complexity:} $O(n)$ time (visits each node once).
\end{itemize}

\subsubsection{Problem: Lowest Common Ancestor of a Binary Tree (Blind 75)}
\begin{itemize}
\item \textbf{Problem:} Find the LCA of two nodes $p$ and $q$ in a general binary tree.
\item \textbf{Solution Idea:} Use a recursive DFS find(node, p, q).
\item \textbf{Base Case:} if node == NULL or node == p or node == q, return node.
\item \textbf{Recurse:} left = find(node.left), right = find(node.right).
\item \textbf{Combine:} if left \&\& right: return node (it's the LCA). else: return left or right (whichever is not NULL).
\item \textbf{Complexity:} $O(n)$ time, $O(h)$ space.
\end{itemize}

\subsubsection{Problem: Construct Tree from Preorder and Inorder (Blind 75)}
\begin{itemize}
\item \textbf{Problem:} Given preorder and inorder traversal arrays, construct the binary tree.
\item \textbf{Solution Idea:} (1) First element of preorder is the root. (2) Find root in inorder to split it into left/right subarrays. (3) Use lengths of these subarrays to split preorder. (4) Recurse.
\item \textbf{Key:} Use a hash map for inorder indices to get $O(1)$ lookup.
\item \textbf{Complexity:} $O(n)$ time, $O(n)$ space.
\end{itemize}

\subsubsection{Problem: Binary Tree Maximum Path Sum (Hard, Blind 75)}
\begin{itemize}
\item \textbf{Problem:} Find the maximum sum of any path in a binary tree.
\item \textbf{Solution Idea:} Use a recursive DFS max\_gain(node). It returns the max path sum \textit{extendable upwards}. It also updates a global max\_sum with the max path sum \textit{contained at this node} (node.key + left\_gain + right\_gain).
\item \textbf{Complexity:} $O(n)$ time, $O(h)$ space.
\end{itemize}

\subsubsection{Problem: AVL Tree Insertion (from GFG Advanced-Data-Structure.pdf)}
\begin{itemize}
\item \textbf{Problem:} Insert a node into an AVL tree and rebalance (Lec 27-29).
\item \textbf{Solution Idea:} (1) Perform standard BST insert. (2) Walk back up from the new node to the root. (3) At each ancestor z, update its height. (4) Find the \textbf{first} ancestor z that violates the AVL property ($|h\_L - h\_R| > 1$). (5) Perform the correct \textbf{Rotation} (LL, RR, LR, RL) at z.
\item \textbf{Key:} Only \textbf{one} rotation (single or double) is needed to fix an insertion.
\item \textbf{Complexity:} $O(\log n)$ time.
\end{itemize}

\subsubsection{Problem: B-Tree Insertion (from GFG Advanced-Data-Structure.pdf)}
\begin{itemize}
\item \textbf{Problem:} Insert a key into a B-Tree (Lec 32).
\item \textbf{Solution Idea (Pre-emptive Split):} We \textit{never} insert into a full node.
\begin{enumerate}
\item Start at the root. As we traverse \textit{down} to find the correct leaf:
\item \textbf{If we encounter a full node x:} We \textbf{split x first} (promoting its median key to its parent) before moving down into the correct (now non-full) child.
\item \textbf{Finally:} We reach the target leaf, which is \textit{guaranteed} not to be full. Insert key $k$.
\end{enumerate}
\item \textbf{Complexity:} $O(\log\_t n)$ disk accesses.
\end{itemize}

\subsubsection{Problem: Implement Trie (Prefix Tree) (Blind 75)}
\begin{itemize}
\item \textbf{Problem:} Implement a Trie with insert, search, and startsWith.
\item \textbf{Solution Idea:} Use a TrieNode class with children[26] and isEndOfWord.
\item \textbf{Complexity:} All ops are $O(L)$ time, where $L$ is length of the word.
\end{itemize}

\subsubsection{Problem: Auto-complete feature using Trie (from GFG Advanced-Data-Structure.pdf)}
\begin{itemize}
\item \textbf{Problem:} Given a query string s, return all words in the dictionary that have s as a prefix.
\item \textbf{Solution Idea:} (1) search the Trie for the node corresponding to the prefix s. (2) If this node p exists, run a DFS from p to find all words in the subtree rooted at p.
\item \textbf{Complexity:} $O(L + K \cdot M)$ where $L$=length of prefix, $K$=number of suggestions, $M$=avg word length.
\end{itemize}

\subsection{Module 7: Graph Algorithms}

\subsubsection{Problem: Number of Islands (Blind 75)}
\begin{itemize}
\item \textbf{Problem:} Given a 2D grid of '1's (land) and '0's (water), count the islands.
\item \textbf{Solution Idea:} Iterate the grid. If grid[r][c] == '1' and unvisited: increment island\_count, and start a \textbf{BFS or DFS} to visit and mark all connected '1's.
\item \textbf{Complexity:} $O(N \cdot M)$ time.
\end{itemize}

\subsubsection{Problem: Topological Sort (from GFG Graphs.pdf)}
\begin{itemize}
\item \textbf{Problem:} Find a linear ordering of vertices in a DAG such that for every edge $(u, v)$, $u$ comes before $v$ (Lec 38).
\item \textbf{Solution Idea (DFS):} (1) Run DFS on the graph. (2) As each vertex \textbf{finishes} (is marked BLACK), insert it onto the \textbf{front} of a linked list. (3) This list is the topological sort.
\item \textbf{Correctness:} For an edge $(u,v)$, $v$ must finish before $u$ (otherwise it's a back edge/cycle). Since $u$ finishes later, it gets added to the front of the list, placing it before $v$.
\item \textbf{Complexity:} $O(n+m)$ time.
\end{itemize}

\subsubsection{Problem: Course Schedule (Lec 38, Blind 75)}
\begin{itemize}
\item \textbf{Problem:} Can you finish all courses given prerequisites?
\item \textbf{Solution Idea:} This is \textbf{Cycle Detection in a Directed Graph}.
\item \textbf{Algorithm (DFS):} Use 3 "colors": WHITE (unvisited), GRAY (in current recursion stack), BLACK (finished). If DFS visits a GRAY node, a \textbf{back edge} (cycle) exists.
\item \textbf{Complexity:} $O(n+m)$ time.
\end{itemize}

\subsubsection{Problem: Detect Cycle in Undirected Graph (from GFG Graphs.pdf)}
\begin{itemize}
\item \textbf{Problem:} Check if an undirected graph has a cycle.
\item \textbf{Solution Idea 1 (DFS):} Run DFS. Keep track of each node's parent. If you find an edge to an already-visited node v that is \textbf{not} your parent, you have found a back edge/cycle.
\item \textbf{Solution Idea 2 (DSU):} (Lec 41) Iterate all edges $(u, v)$. For each edge, if Find(u) == Find(v): a cycle exists (they are already connected). else: Union(u, v).
\item \textbf{Complexity:} $O(n+m)$ for DFS. $O(m \cdot \alpha(n))$ for DSU.
\end{itemize}

\subsubsection{Problem: Snake and Ladder (from GFG Graphs.pdf)}
\begin{itemize}
\item \textbf{Problem:} Find the minimum number of dice throws to win a Snake and Ladder game.
\item \textbf{Solution Idea:} This is a \textbf{Shortest Path in an Unweighted Graph} problem.
\item \textbf{Graph:} The board (squares 1-100) is the set of $n$ vertices.
\item \textbf{Edges:} From each square i, there are 6 edges to i+1, ..., i+6 (dice rolls).
\item \textbf{Snakes/Ladders:} These are not edges, but "teleports". If a ladder is at i to j, you only have an edge to j, not i. (Model this in the adj. list).
\item \textbf{Algorithm:} Run \textbf{BFS} from square 1. The dist[100] is the answer.
\item \textbf{Complexity:} $O(n)$ time (where $n=100$).
\end{itemize}

\subsubsection{Problem: Word Search (Backtracking/DFS)}
\begin{itemize}
\item \textbf{Problem:} Find if a word exists in a 2D grid. Cannot reuse a cell.
\item \textbf{Solution Idea:} \textbf{DFS with backtracking}. Start DFS from any cell matching word[0]. Mark cell as visited, explore 4 neighbors, then un-mark cell.
\item \textbf{Complexity:} $O(N \cdot M \cdot 3^L)$ where $L$ is word length.
\end{itemize}

\subsubsection{Problem: Kruskal's Algorithm for MST (from GFG Graphs.pdf)}
\begin{itemize}
\item \textbf{Problem:} Find a Minimum Spanning Tree (Lec 41).
\item \textbf{Solution Idea:} (1) \textbf{Sort} all $m$ edges by weight ($O(m \log m)$). (2) Initialize a \textbf{DSU} with $n$ sets. (3) Iterate sorted edges $(u, v)$: if Find(u) != Find(v): add $(u, v)$ to MST and Union(u, v).
\item \textbf{Complexity:} $O(m \log m)$ or $O(m \log n)$ (dominated by sort).
\end{itemize}
% --- MODULE 9: Exercise & Challenge Solutions ---
% --- Sourced from Lecture Slides & CLRS ---

\section*{Module 9: Exercise \& Challenge Solutions}
This module contains the solutions to the exercises and challenges presented in the lecture slides.

\subsection*{Module 2: Stacks \& Queues}

\subsubsection*{Problem: Efficient Array-Based Queue (Lec 8, Slide 54)}
\begin{itemize}
    \item \textbf{Problem:} Provide an efficient array-based implementation of a queue that requires $O(n)$ space (where $n$ is max size) but $n$ is not known beforehand.
    \item \textbf{Solution Idea:} Combine two techniques from the lectures.
    \begin{enumerate}
        \item \textbf{Circular Array (Lec 8):} Use `front` and `back` indices with the modulo operator (`\% N`, where $N$ is array capacity). This makes `Dequeue` an $O(1)$ operation by just moving the `front` pointer, avoiding an $O(N)$ shift.
        \item \textbf{Doubling/Halving (Lec 9):} To handle unknown size, start with a small array.
        \begin{itemize}
            \item \textbf{On `Enqueue`:} If the array is full, create a new array of double the size ($2N$) and copy the $N$ elements over. This gives an amortized $O(1)$ cost.
            \item \textbf{On `Dequeue`:} If the number of elements falls below $N/4$, create a new array of half the size ($N/2$) and copy the elements. This prevents "thrashing" and maintains the amortized $O(1)$ cost.
        \end{itemize}
    \end{enumerate}
    \item \textbf{Result:} All operations (`Enqueue`, `Dequeue`) achieve $\mathbf{O(1)}$ amortized time complexity.
\end{itemize}

\subsection*{Module 3: Trees \& Traversals}

\subsubsection*{Problem: Traversal via Stacks and Queues (Lec 14, Slide 31)}
\begin{itemize}
    \item \textbf{Level-Order (using a Queue):} This is the BFS algorithm.
    \begin{enumerate}
        \item Create a Queue, `q.enqueue(root)`.
        \item While $q$ is not empty: $u = q.dequeue()$.
        \item Visit $u$.
        \item If `u.left != NULL`, `q.enqueue(u.left)`.
        \item If `u.right != NULL`, `q.enqueue(u.right)`.
    \end{enumerate}
    \item \textbf{Preorder (using a Stack):}
    \begin{enumerate}
        \item Create a Stack, `s.push(root)`.
        \item While $s$ is not empty: $u = s.pop()$.
        \item Visit $u$.
        \item If `u.right != NULL`, `s.push(u.right)`. (Push right first)
        \item If `u.left != NULL`, `s.push(u.left)`. (Push left last)
    \end{enumerate}
    \item \textbf{Inorder (using a Stack):} This is the most complex iterative traversal.
    \begin{enumerate}
        \item Create a Stack $s$. `curr = root`.
        \item While `curr != NULL` or $s$ is not empty:
        \item \textbf{Go left:} While `curr != NULL`: `s.push(curr)`, `curr = curr.left`.
        \item \textbf{Visit:} `curr = s.pop()`. Visit `curr`.
        \item \textbf{Go right:} `curr = curr.right`.
    \end{enumerate}
\end{itemize}

\subsubsection*{Problem: Count Trees from Level-Order (Lec 14, Slide 32)}
\begin{itemize}
    \item \textbf{Challenge:} "How many binary trees have the sequence ABCDEFG as their level-order traversal?"
    \item \textbf{Solution:} The level-order traversal sequence does \textit{not} uniquely define a binary tree, because it does not encode `NULL` children.
    \item \textbf{Example 1:} A \textit{full, complete} tree: `A.left=B`, `A.right=C`, `B.left=D`, `B.right=E`, `C.left=F`, `C.right=G`. This produces `A B C D E F G`.
    \item \textbf{Example 2:} A "path" tree: `A.right=B`, `B.right=C`, `C.right=D`, `D.right=E`, `E.right=F`, `F.right=G`. This also produces `A B C D E F G`.
    \item \textbf{Conclusion:} There are many possible trees. However, as noted on Lec 14, Slide 47, if the tree is constrained to be an \textbf{Almost Complete Binary Tree}, the sequence uniquely defines \textbf{exactly one} tree.
\end{itemize}

\subsection*{Module 6: BSTs \& AVL Trees}

\subsubsection*{Problem: Inorder Traversal of BST is Sorted (Lec 25, Slide 6)}
\begin{itemize}
    \item \textbf{Proof (by Induction):}
    \item \textbf{Base Case:} $n=0$ (empty tree) or $n=1$ (root only). The traversal is empty or has one node, which is sorted.
    \item \textbf{Inductive Hypothesis (IH):} Assume that for any BST with $k < n$ nodes, an inorder traversal produces a sorted list.
    \item \textbf{Inductive Step (for $n$ nodes):} Let $T$ be a BST with $n$ nodes and root $r$. $T$'s left subtree $T\_L$ and right subtree $T\_R$ are both BSTs with $< n$ nodes.
    \begin{enumerate}
        \item By the \textbf{BST Invariant} (Lec 24), all keys in $T\_L$ are $\le r.\text{key}$, and all keys in $T\_R$ are $\ge r.\text{key}$.
        \item An inorder traversal of $T$ is: `Inorder(T\_L)`, then `r.key`, then `Inorder(T\_R)`.
        \item By the (IH), `Inorder(T\_L)` is a sorted list.
        \item By the (IH), `Inorder(T\_R)` is a sorted list.
        \item Combining these facts, the full traversal is: [sorted list $\le r$] + [$r$] + [sorted list $\ge r$].
        \item This final combined list is fully sorted.
    \end{enumerate}
\end{itemize}

\subsubsection*{Problem: Find All Occurrences of Key k (Lec 25, Slide 17)}
\begin{itemize}
    \item \textbf{Problem:} Design an algorithm to find all occurrences of $k$.
    \item \textbf{Solution Idea:} The standard `Search` stops at the first $k$. Since duplicates $\le k$ can be in the left subtree and $\ge k$ can be in the right, we must search both.
    \item \textbf{Algorithm} `FindAll(node, k, results\_list)`:
    \begin{enumerate}
        \item If `node == NULL`, return.
        \item If `k < node.key`:
        \item `FindAll(node.left, k, results\_list)`
        \item Else if `k > node.key`:
        \item `FindAll(node.right, k, results\_list)`
        \item Else (`k == node.key`):
        \item `results\_list.add(node)`
        \item `FindAll(node.left, k, results\_list)` (to find other $k$'s in left)
        \item `FindAll(node.right, k, results\_list)` (to find other $k$'s in right)
    \end{enumerate}
    \item \textbf{Complexity:} $O(h + c)$, where $h$ is tree height and $c$ is the number of occurrences.
\end{itemize}

\subsubsection*{Problem: Compute Predecessor (Lec 25, Slide 32)}
\begin{itemize}
    \item \textbf{Problem:} Modify the `Successor` algorithm to compute `Predecessor(x)`.
    \item \textbf{Solution Idea:} This is the symmetric (mirror-image) of `Successor`.
    \item \textbf{Algorithm:}
    \begin{enumerate}
        \item \textbf{Case 1: If $x$ has a left subtree.}
        The predecessor is the \texttt{Maximum()} node in $x$'s left subtree.
        \item \textbf{Case 2: If $x$ has NO left subtree.}
        Go up the tree using `parent` pointers. The predecessor is the \textbf{first ancestor} $y$ such that $x$ is in the \textbf{right subtree} of $y$.
    \end{enumerate}
    \item \textbf{Complexity:} $O(h)$.
\end{itemize}

\subsubsection*{Problem: Expected Height of Random BST (Lec 25, Slide 52)}
\begin{itemize}
    \item \textbf{Problem:} What is the expected height of a BST constructed from a random permutation of $\{1, \dots, n\}$?
    \item \textbf{Solution (from CLRS 12.4):}
    \item \textbf{Result:} The expected height is $\mathbf{O(\log n)}$.
    \item \textbf{Proof Sketch:} Let $X\_n$ be the height of a random BST on $n$ nodes. Let $R$ be the rank of the key chosen as the root (with $P(R=i) = 1/n$).
    \item The height of the tree is $1 + \max(X\_{i-1}, X\_{n-i})$, where $X\_{i-1}$ and $X\_{n-i}$ are the heights of the left/right subtrees.
    \item The analysis of this recurrence (which is very similar to the analysis of randomized quicksort) shows that $E[X\_n] = O(\log n)$.
\end{itemize}

\subsubsection*{Problem: Maintain Subtree Size in BST (Lec 26, Slide 5)}
\begin{itemize}
    \item \textbf{Problem:} Update `Insert` and `Delete` to maintain the `size` field at every node.
    \item \textbf{Solution (`size(x) = size(x.left) + size(x.right) + 1`):}
    \item \textbf{Insert(z):}
    \begin{enumerate}
        \item Perform a normal BST insert. `size(z)` is initialized to 1.
        \item As the search path is traversed *down* to find the insertion spot, increment the `size` of every node on that path.
    \end{enumerate}
    \item \textbf{Delete(z):}
    \begin{enumerate}
        \item Perform a normal BST delete. Let $y$ be the node that was physically removed (either $z$ or $z$'s successor).
        \item Starting from `parent(y)`, walk *up* the tree to the root, decrementing the `size` of every node on that path.
    \end{enumerate}
    \item \textbf{Complexity:} Both operations remain $O(h)$.
\end{itemize}

\subsubsection*{Problem: AVL Rotation Height Proof (Lec 28, Slides 42, 47)}
\begin{itemize}
    \item \textbf{Problem:} In Case 1(a) (Right-Rotate) and Case 2(a) (Left-Rotate), prove the middle subtree (yellow box) must have height $h-3$.
    \item \textbf{Proof (for Case 1a / Right-Rotate):}
    \begin{enumerate}
        \item Let $Y$ be the unbalanced node. Its height \textit{after} insertion is $h$.
        \item Let $X$ be its left child. Since this is Case 1a, $X$ is left-heavy and $Y$ is unbalanced.
        \item This means `h(X) = h-1` and `h(Y.right) = h-3`.
        \item The imbalance was caused by an insertion into $X$'s left subtree, so $X$'s height just grew from $h-2$ to $h-1$.
        \item Since $Y$ *was* balanced before the insert (when `h(X) = h-2`), its other child `Y.right` must have had height $h-3$ (as $h-2$ would also be balanced).
        \item Since $X$ *was* balanced before the insert (at height $h-2$), and the insert occurred in its \textit{left} child (making it $h-2$), its \textit{right} child (the yellow box) must have had height $\le h-3$.
        \item To maintain $X$'s balance *before* the insert, if $X.left$ was $h-3$, $X.right$ must have been $h-3$.
        \item Thus, the yellow box's height *before* insertion was $h-3$, and it was unchanged by the insertion.
    \end{enumerate}
\end{itemize}

% --- END OF CONTENT ---

\end{document}